% =====================================================================
%  IEEE Conference Template（支援繁體中文）
%  ---------------------------------------------------------------------
%  編譯器：必須使用 XeLaTeX
%    Overleaf → 左上角 Menu → Settings → Compiler → XeLaTeX
%  參考文獻：寫在 refs.bib，用 \cite{key} 引用，格式由 IEEEtran.bst 產生
%  圖片：放在 figures/ 資料夾，再把 example-image-a 換成你的檔名
% =====================================================================
\RequirePackage{iftex}
\RequireXeTeX  % 沒切成 XeLaTeX 時，會直接報錯提醒你

\documentclass[conference]{IEEEtran}
% 期刊格式改成：\documentclass[journal]{IEEEtran}

% ---------- 字型（英文 Times、中文思源宋體／黑體） ----------
\usepackage{fontspec}
\setmainfont{TeX Gyre Termes}        % Times New Roman 的開源替代
\usepackage{xeCJK}
\setCJKmainfont{Noto Serif CJK TC}   % 中文內文：思源宋體 繁中
\setCJKsansfont{Noto Sans CJK TC}    % 中文無襯線：思源黑體 繁中
\setCJKmonofont{Noto Sans Mono CJK TC}
\xeCJKsetup{CJKecglue={\hskip 0.15em plus 0.05em minus 0.05em}} % 中英文之間的間距

% ---------- 常用套件 ----------
\usepackage{cite}                    % IEEE 風格的引用編號 [1]–[3]
\usepackage{amsmath,amssymb,amsfonts}
\usepackage{algorithmic}             % IEEE 建議的演算法環境
\usepackage{algorithm}
\usepackage{graphicx}
\graphicspath{{figures/}}
\usepackage{booktabs}                % 三線表 \toprule \midrule \bottomrule
\usepackage{multirow}
\usepackage{textcomp}
\usepackage{xcolor}
\usepackage{url}
\usepackage[hidelinks]{hyperref}     % 可點擊的連結與交叉引用

% 演算法標題改成中文的話取消註解：
% \floatname{algorithm}{演算法}
% \renewcommand{\algorithmicrequire}{\textbf{輸入：}}
% \renewcommand{\algorithmicensure}{\textbf{輸出：}}

\begin{document}

% =====================================================================
%  標題與作者
% =====================================================================
\title{論文標題：以物聯網感測資料進行○○分析與預測\\
{\large Paper Title: IoT Sensor Data Analysis and Prediction for ○○}}

\author{
\IEEEauthorblockN{李偉聖}
\IEEEauthorblockA{\textit{○○學系} \\
\textit{國立中興大學}\\
臺中市, 臺灣 \\
email@example.com}
\and
\IEEEauthorblockN{第二作者}
\IEEEauthorblockA{\textit{○○學系} \\
\textit{國立中興大學}\\
臺中市, 臺灣 \\
email@example.com}
}

\maketitle

% =====================================================================
%  摘要與關鍵字
% =====================================================================
\begin{abstract}
本文在此簡述研究動機、使用的方法、資料來源與主要結果。摘要建議控制在
150 到 250 字之間，不放引用、圖表或公式。例如：本研究蒐集○○感測器的
溫溼度資料，使用○○模型進行預測，實驗結果顯示 RMSE 為 ○○，比基準方法
降低 ○○\%。
\end{abstract}

\begin{IEEEkeywords}
物聯網, 感測資料分析, 機器學習, 時間序列預測
\end{IEEEkeywords}

% =====================================================================
\section{Introduction（緒論）}
\label{sec:intro}
% =====================================================================
說明研究背景與問題。物聯網（Internet of Things, IoT）讓大量裝置能夠互相連線
並交換資料~\cite{atzori2010iot,gubbi2013iot}。接著寫出目前遇到的問題，以及
本文想解決什麼。

本文的主要貢獻如下：
\begin{itemize}
  \item 貢獻一：例如建置一套○○資料蒐集系統。
  \item 貢獻二：例如提出○○分析方法。
  \item 貢獻三：例如以實際資料驗證方法的效果。
\end{itemize}

本文其餘章節安排如下：第~\ref{sec:related}~節回顧相關研究，
第~\ref{sec:system}~節介紹系統架構，第~\ref{sec:method}~節說明方法，
第~\ref{sec:exp}~節呈現實驗結果，第~\ref{sec:conclusion}~節為結論。

% =====================================================================
\section{Related Work（相關研究）}
\label{sec:related}
% =====================================================================
整理過去的研究，並說明本文與它們的差異。例如 Random Forest~\cite{breiman2001rf}
常用於感測資料分類，實作上可使用 scikit-learn~\cite{pedregosa2011sklearn}。

% =====================================================================
\section{System Architecture（系統架構）}
\label{sec:system}
% =====================================================================
描述整體系統，例如感測層、網路層與應用層。圖~\ref{fig:arch}~為系統架構圖。

\begin{figure}[htbp]
  \centering
  \includegraphics[width=\columnwidth]{example-image-a}
  % 換成自己的圖：把檔案放進 figures/，例如 \includegraphics[width=\columnwidth]{architecture.png}
  \caption{系統架構圖（單欄寬度）。}
  \label{fig:arch}
\end{figure}

% =====================================================================
\section{Methodology（研究方法）}
\label{sec:method}
% =====================================================================
\subsection{資料前處理}
說明資料清理、缺值處理與正規化。例如 Min-Max 正規化如式~\eqref{eq:minmax}：
\begin{equation}
  x' = \frac{x - x_{\min}}{x_{\max} - x_{\min}}
  \label{eq:minmax}
\end{equation}

\subsection{模型與評估指標}
模型預測誤差以均方根誤差（RMSE）衡量：
\begin{equation}
  \mathrm{RMSE} = \sqrt{\frac{1}{N}\sum_{i=1}^{N}\left(y_i - \hat{y}_i\right)^2}
  \label{eq:rmse}
\end{equation}
其中 $y_i$ 為實際值，$\hat{y}_i$ 為預測值，$N$ 為樣本數。

\subsection{演算法流程}
完整流程如 Algorithm~\ref{alg:flow} 所示。

\begin{algorithm}[htbp]
  \caption{資料分析流程}
  \label{alg:flow}
  \begin{algorithmic}[1]
    \REQUIRE 感測資料集 $D$
    \ENSURE 訓練好的模型 $M$
    \STATE 移除缺值與異常值
    \STATE 依式~\eqref{eq:minmax} 正規化 $D$
    \STATE 將 $D$ 切分為訓練集 $D_{train}$ 與測試集 $D_{test}$
    \FOR{每組超參數 $\theta$}
      \STATE 以 $D_{train}$ 訓練模型 $M_\theta$
      \STATE 計算 $M_\theta$ 在驗證集上的 RMSE
    \ENDFOR
    \RETURN RMSE 最低的模型 $M$
  \end{algorithmic}
\end{algorithm}

% =====================================================================
\section{Experiments and Results（實驗與結果）}
\label{sec:exp}
% =====================================================================
\subsection{實驗設定}
說明資料來源、資料筆數、硬體環境與參數設定。

\subsection{實驗結果}
表~\ref{tab:result}~比較各方法的表現。

\begin{table}[htbp]
  \centering
  \caption{各模型預測結果比較}
  \label{tab:result}
  \begin{tabular}{lccc}
    \toprule
    模型 & RMSE $\downarrow$ & MAE $\downarrow$ & $R^2$ $\uparrow$ \\
    \midrule
    Linear Regression & 0.00 & 0.00 & 0.00 \\
    Random Forest     & 0.00 & 0.00 & 0.00 \\
    本文方法          & \textbf{0.00} & \textbf{0.00} & \textbf{0.00} \\
    \bottomrule
  \end{tabular}
\end{table}
% 注意：IEEE 格式的表格標題放在表格「上方」，圖片標題放在圖片「下方」。

圖~\ref{fig:result}~為跨雙欄的結果圖。

\begin{figure*}[htbp]
  \centering
  \includegraphics[width=0.48\textwidth]{example-image-b}
  \hfill
  \includegraphics[width=0.48\textwidth]{example-image-c}
  \caption{預測結果比較（跨雙欄寬度）：(左) 實際值與預測值；(右) 誤差分布。}
  \label{fig:result}
\end{figure*}

% =====================================================================
\section{Discussion（討論）}
% =====================================================================
分析結果背後的原因、方法的限制，以及可能的改進方向。

% =====================================================================
\section{Conclusion（結論）}
\label{sec:conclusion}
% =====================================================================
總結本文的做法與主要發現，並提出未來工作。

% =====================================================================
\section*{Acknowledgment（致謝）}
% =====================================================================
感謝○○提供的資料與協助。

% =====================================================================
%  參考文獻（內容寫在 refs.bib）
% =====================================================================
\bibliographystyle{IEEEtran}
\bibliography{refs}

\end{document}
